\documentclass[12pt]{article}
\usepackage{amsmath,amssymb}

% -- Page Margin --
\usepackage[margin=1in]{geometry}


\title
{
    Teoría de la Computación 2026-1 \\
    Tarea 2
}

    \begin{document}

    \maketitle

    \begin{center}
        \begin{tabular}{r|l}
            \textbf{Expediente} & \textbf{Nombre} \\ \hline
            223209096 & Ahumada Herrera Jorge Alán \\
            219208106 & Bórquez Guerrero Angel Fernando \\
            222206586 & Juan Valenzuela Gael \\
            223201053 & Solis Zatarain Owen Adiel \\
        \end{tabular}
    \end{center}

    \rule{\linewidth}{0.3mm}

    \vspace{0.3cm}



\textbf{Ejercicio 16.} Para cada $n > 0$ fijo, construya un AFD de $n$ estados que distinga solamente a las cadenas de $\{a\}^*$ de la forma $a^{kn}$ donde $k \ge 0$.

\vspace{1em}

\begin{flushleft}
$A = (Q, \Sigma, \delta, q_0, F)$\\[1em]

$Q = \{q_0, q_1, q_2, \dots, q_{n-1}\}$\\[1em]
    
$\Sigma = \{a\}$\\[1em]
    
$q_0 = q_0$\\[1em]
    
$F = \{q_0\}$\\[1em]
    
$\delta(q_i, a) = q_{(i+1) \bmod n} \quad \text{para todo } 0 \le i \le n-1$
\end{flushleft}



\newpage
\textbf{Ejercicio 17.} Considere una máquina despachadora de bebidas. 
Esta máquina solo vende dos tipos de bebida: agua purificada y horchata. El agua vale 15
pesos y la horchata vale 20 pesos. La máquina solo acepta monedas de 5 y de
10 pesos y no da cambio. Diseñe un AFD que modele el conteo de las monedas
que el cliente está insertando. Ya que la máquina no da cambio, los estados
de aceptación serán aquellos en el que el cliente haya insertado una cantidad
suficiente para comprar n aguas y m horchatas (por supuesto m ó n pueden
ser cero, pero no ambas a la vez, es decir si el cliente no inserta nada, no se
estaría en un estado de aceptación). Por ejemplo, si el cliente introduce en total
25 pesos no se estaría en un estado de aceptación, ya que se podría comprar
una horchata pero sobrarían 5 pesos. Si el cliente inserta en total 55 pesos, sí se
estaría en un estado de aceptación ya que se podrían comprar 2 horchatas y un
agua purificada exactamente. \\

\( AFD = \{ Q, \Sigma, \delta, q_{0}, F \} \) \\

\( Q = \{ q_{0}, 5, 10, 15, 20, 25, 30 \} \) \\ 

\( \Sigma = \{ 5, 10 \} \) \\

\( F = \{ 15, 20, 30 \} \) \\

\begin{tabular}{c|cc}
    $\delta$ & 5 & 10 \\
    \hline
    $q_{0}$ & 5 & 10 \\
    5 & 10 & 15 \\
    10 & 15 & 20 \\
    15 & 20 & 25 \\
    20 & 25 & 30 \\
    25 & 30 & 30 \\
    30 & 30 & 30 

\end{tabular} \\



\newpage
\textbf{Ejercicio 19.} Sea $A=(Q,\Sigma,\delta,q_{0},F)$ un AFD y sea $A^{c}=(Q,\Sigma,\delta,q_{0},Q\setminus F)$. Pruebe que $L(A^{c})=L(A)^{c}$.

\vspace{1.5em}

\textbf{Demostraci\'on:}

Sea $w \in \Sigma^*$ una cadena cualquiera.

\begin{flalign*}
w \in L(A^c) &\iff \hat{\delta}(q_0, w) \in (Q \setminus F) &\\[1em]
&\iff \hat{\delta}(q_0, w) \in Q \quad \text{y} \quad \hat{\delta}(q_0, w) \notin F &\\[1em]
&\iff \hat{\delta}(q_0, w) \notin F \quad \text{(dado que } \hat{\delta}(q_0, w) \in Q \text{ siempre en un AFD)} &\\[1em]
&\iff w \notin L(A) &\\[1em]
&\iff w \in L(A)^c &
\end{flalign*}

Por lo tanto, se concluye que $L(A^c) = L(A)^c$. \hfill $\blacksquare$



\textbf{Ejercicio 20.} Sea $A=(Q,\Sigma,\delta,q_0,F)$ el AFD de la figura, donde:
\[
Q=\{q_0,q_1,q_2\}, \qquad 
\Sigma=\{0,1\},
\]
\[
F=\{q_0,q_2\}.
\]

Las transiciones son:

\[
\begin{aligned}
\delta(q_0,0)&=q_1, & \delta(q_0,1)&=q_2, \\
\delta(q_1,0)&=q_1, & \delta(q_1,1)&=q_2, \\
\delta(q_2,0)&=q_1, & \delta(q_2,1)&=q_2.
\end{aligned}
\]

Queremos demostrar que:

\[
L(A)=\{\varepsilon\}\cup\{x1 : x\in\{0,1\}^*\}.
\]

Sea 
\[
S=\{\varepsilon\}\cup\{x1 : x\in\{0,1\}^*\}.
\]

Demostraremos que $L(A)=S$ probando doble inclusión.

\paragraph{Nota.}
Para demostrar que $A=B$ se usa \textbf{doble inclusión}: 
se prueba que $A\subseteq B$ y que $B\subseteq A$. 
Si ambos contienen exactamente los mismos elementos, entonces son iguales.


\newpage
\subsection*{1. $L(A)\subseteq S$}

Sea $w\in L(A)$. Entonces $\hat{\delta}(q_0,w)\in F$.

Comportamiento del autómata:

\begin{itemize}
\item $q_0$ es estado inicial y de aceptación.
\item $q_2$ es estado de aceptación.
\item $q_1$ no es de aceptación.
\end{itemize}

Observamos que:

\begin{itemize}
\item Cada vez que se lee un símbolo $1$, el autómata pasa a $q_2$.
\item Cada vez que se lee un símbolo $0$, el autómata pasa a $q_1$.
\end{itemize}

Por lo que, el estado en el que termina el autómata depende únicamente del \textbf{último símbolo} de la cadena:

\begin{itemize}
\item Si la cadena termina en $1$, termina en $q_2$ (estado de aceptación).
\item Si la cadena termina en $0$, termina en $q_1$ (no aceptación).
\end{itemize}

Aparte, la cadena vacía $\varepsilon$ es aceptada porque el estado inicial $q_0$ es de aceptación.

Entonces, si $w\in L(A)$, necesariamente:

\begin{itemize}
\item $w=\varepsilon$, o
\item $w$ termina en $1$.
\end{itemize}

En el segundo caso, existe $x\in\{0,1\}^*$ tal que $w=x1$.

Por lo tanto, $w\in S$. Concluimos que:

\[
L(A)\subseteq S.
\]

\newpage
\subsection*{2. $S\subseteq L(A)$}

Sea ahora $w\in S$.

\begin{itemize}
\item Si $w=\varepsilon$, entonces $\delta^*(q_0,\varepsilon)=q_0\in F$, por lo tanto $\varepsilon\in L(A)$.
\item Si $w=x1$, entonces el último símbolo de $w$ es $1$. Como vimos antes, al leer un $1$ el autómata termina en $q_2$, que es estado de aceptación.
\end{itemize}

Por lo tanto, toda cadena que termina en $1$ es aceptada por el autómata.

Así, $w\in L(A)$ y se cumple:

\[
S\subseteq L(A).
\]

\subsection*{Conclusión}

Probamos que:

\[
L(A)\subseteq S
\quad \text{y} \quad
S\subseteq L(A).
\]

Por lo tanto,

\[
L(A)=\{\varepsilon\}\cup\{x1 : x\in\{0,1\}^*\}.
\]

\hfill $\square$




















\end{document}
